\chapter{Planificación}

A la hora de trabajar en un proyecto de estas características, se hace necesario diseñar un plan de trabajo realista que permita estructurar el esfuerzo, establecer una rutina de progreso y anticipar los riesgos derivados de las limitaciones propias de un Trabajo de Fin de Grado. Por ello, en este capítulo se abordan las cuestiones relativas a la planificación del proyecto, tanto desde un punto de vista metodológico como desde una perspectiva temporal y de esfuerzo, desglosando las fases del trabajo, la estimación de horas y el calendario previsto. Asimismo, se exponen los condicionantes que han influido en la planificación y las decisiones tomadas al respecto.

\section{Metodología de trabajo}
\label{sec:metodologia}

El modelo de desarrollo adoptado para afrontar este proyecto ha sido el desarrollo iterativo. Según la Agile Alliance \parencite{agilealliance_iterative}, los proyectos iterativos se caracterizan por permitir de forma intencionada la repetición de las actividades de desarrollo y la revisión de los mismos productos de trabajo a lo largo del tiempo, lo que se conoce como reelaboración planificada. Además, este enfoque suele estructurarse en torno a una serie de iteraciones de duración fija, si bien su característica principal no es la periodicidad de los ciclos, sino la posibilidad de refinar progresivamente el resultado hasta lograr un proyecto que cumpla con las necesidades del cliente.

Aplicado a este Trabajo de Fin de Grado, el desarrollo iterativo se ha materializado en una sucesión de ciclos en los que el autor ha seguido un plan consistente en una serie de fases: planificación, toma de requisitos, análisis, diseño, implementación y pruebas. Al término de cada iteración, el resultado parcial se ha presentado al tutor del proyecto, que ha asumido el rol de cliente y ha proporcionado retroalimentación sobre el trabajo realizado. Dicha retroalimentación se ha incorporado como punto de partida de la siguiente iteración, permitiendo así corregir desviaciones, reajustar prioridades y refinar el producto de forma continua.

Esta dinámica resulta especialmente adecuada para un proyecto de las características de Liber, en el que el alcance funcional ha ido evolucionando conforme se adquiría un mayor conocimiento del dominio del problema y de las necesidades reales de las bibliotecas a las que se orienta la aplicación. La naturaleza iterativa del proceso ha permitido, además, compatibilizar el desarrollo del software con la redacción de la presente memoria, dado que cada iteración ha generado a la vez avances técnicos y documentación asociada a los mismos.
Se procede ahora a describir los elementos que han marcado la metodología de trabajo utilizada:
\begin{enumerate}
	\item Roles: Al haber seguido una metodología iterativa, no se han adoptado roles más allá del desarrollador de la aplicación (adoptado por el alumno), que se encarga de todo lo relacionado con la construcción del proyecto, y el cliente (rol desempeñado por el tutor), quien evalúa la aplicación y facilita la retroalimentación para acercar la aplicación a sus necesidades.
	\item Eventos: Aunque el desarrollo iterativo no prescribe un calendario de reuniones concreto, en este proyecto se ha establecido un ritmo	de encuentros periódicos con el tutor con una frecuencia aproximada de una a dos reuniones al mes. Estas reuniones han supuesto el principal mecanismo de sincronización y validación del avance del trabajo. En cada encuentro se  solía presentar el estado actual del proyecto, mostrando las nuevas características y funcionalidades implementadas desde la reunión anterior, así como los avances en la redacción de la memoria. A partir de esta revisión conjunta, se discutían tres cuestiones principales: la dirección que estaba tomando el proyecto, las tareas que debían abordarse a continuación una vez revisado lo presentado, y las posibles mejoras o ajustes que podrían adoptarse para futuras iteraciones. De este modo, las reuniones servían como un punto de control del progreso, así como espacio para valorar el punto en que se encontraba el trabajo y tomar decisiones de forma consensuada con el tutor, que asumía el rol de cliente en la validación de cada iteración.
	\item Artefactos: los artefactos generados consisten en un diario de actividad, en el que se refleja toda labor llevada a cabo relacionada con el trabajo, y un libro de actas redactado en un documento de texto en el que se encuentra tanto las órdenes del día de las reuniones periódicas, como el resultado y las decisiones tomadas en la misma.
	\item Entorno de trabajo: El canal seleccionado junto a la periodicidad de las reuniones debía facilitar la coordinación entre alumno y tutor. De esta forma se decidió tener entre una y dos semanas periódicas como se ha mencionado, y se tomó la decisión de asistir (en la medida de lo posible) presencialmente. Como no siempre pudo ser así, también han ocurrido reuniones vía videollamada en algunas ocasiones, lo cual no influía en el curso de la reunión. Adicionalmente, ha habido también comunicación vía correo electrónico para dudas urgentes puntuales, entrega de documentos como partes de la memoria para su corrección, o acordar día y medio por el que tener la reunión. Ha sido una comunicación constante que ocurría todas las semanas, facilitando que el tutor pudiera corregir o retroalimentar constantemente.
\end{enumerate}
En síntesis, la metodología iterativa aplicada se ha articulado en torno a dos roles claramente diferenciados (desarrollador y cliente), un ritmo constante de reuniones periódicas con el tutor, un conjunto sencillo de artefactos de seguimiento (diario de actividad y libro de actas) y un entorno de trabajo flexible que ha combinado encuentros presenciales, videollamadas y comunicación por correo electrónico. Esta combinación ha permitido mantener un flujo de retroalimentación continuo a lo largo de todo el proyecto, favoreciendo que el producto final se ajustase de forma progresiva a las necesidades planteadas.

\section{Estimación del esfuerzo}
Para la planificación temporal y de costes existen distintos sistemas de estimación, entre los que destacan dos: la estimación mediante puntos de función y su transformación en líneas de código: el modelo COCOMO II, y la estimación mediante puntos de caso de uso, basada en un análisis previo de los requisitos funcionales de la aplicación.
En este Trabajo de Fin de Grado se ha optado por el segundo enfoque, la estimación mediante puntos de caso de uso. Esta elección se justifica por tres motivos. Primeramente, el capítulo de Análisis de esta memoria ya recoge un conjunto de criterios de aceptación asociados a cada objetivo general definido, proporcionando una base sólida y verificable para identificar actores y casos de uso sin necesidad de un esfuerzo adicional de modelado. En segundo lugar, el método de puntos de caso de uso no requiere disponer de datos históricos de proyectos anteriores, a diferencia de COCOMO II, cuyos resultados dependen en gran medida de la disponibilidad de métricas históricas representativas del entorno de desarrollo, artefacto del cual no se dispone. En tercer lugar, al ser un proyecto individual y estar acotado su alcance, la simplicidad del método y su compatibilidad con el análisis funcional ya realizado reducen la posibilidad de introducir un sobreesfuerzo de estimación desproporcionado respecto al tamaño del trabajo. Adicionalmente, resulta especialmente adecuado para la metodología de desarrollo elegida; tratándose de un desarrollo iterativo, se revisan constantemente los casos de uso a elaborar, dotando de un mayor control y revisión sobre los mismos y por tanto de su estimación.

Sin embargo, no hay que olvidar el sentido de esta herramienta; su propio nombre dice que es una estimación, por lo que se remarca su carácter provisional y sujeto a revisión conforme el desarrollo avanza y el conocimiento del dominio se consolida. En las secciones siguientes se detalla la aplicación concreta del método elegido y los resultados obtenidos.
\section{Planificación temporal}
Este TFG comenzó a finales de febrero de 2026, en principio iba a terminar en mayo o junio, pero por cuestiones externas ha terminado en noviembre. La carga de trabajo estimada es de 300 horas, pero habiendo considerado los bloqueos y parones, ha terminado siendo un trabajo de 400 horas.
\subsection{Planificación temporal inicial}
\subsection{Seguimiento del proyecto}
\section{Presupuesto económico}
Evaluar el presupuesto usando métodos conocidos y luego hacer una comparación entre ellos.
\subsection{Comparativa entre modelos de estimación}
\section{Balance final}